% ============================================================
% A Finite-Cutoff Hilbert-Space Model for Prime--Zero
% Energy Structure
%
% Author:   Ulrich Tehrani
% Date:     May 2026
% Version: v14
% Zenodo:   https://doi.org/10.5281/zenodo.19307989
% License:  CC BY 4.0
% ============================================================

\documentclass[11pt]{article}
\usepackage{cmap}            % ToUnicode CMaps for §, en/em-dashes etc.
\usepackage[T1]{fontenc}
\emergencystretch=2em        % avoid overfull hboxes from inline math
\usepackage[utf8]{inputenc}
\usepackage[a4paper,margin=1in]{geometry}
\usepackage{amsmath,amssymb,amsthm,mathtools}
\usepackage{xcolor}
\usepackage{hyperref}
\usepackage{enumitem}
\usepackage{setspace}
\usepackage{booktabs}
\usepackage{array}
\usepackage{graphicx}

\hypersetup{colorlinks=true,linkcolor=blue!50!black,urlcolor=blue!50!black,citecolor=blue!50!black}
\setlength{\parskip}{0.5em}
\setlength{\parindent}{0pt}
\onehalfspacing

% -- Theorem environments 
\newtheorem{theorem}{Theorem}[section]
\newtheorem{proposition}[theorem]{Proposition}
\newtheorem{lemma}[theorem]{Lemma}
\newtheorem{corollary}[theorem]{Corollary}
\newtheorem{definition}[theorem]{Definition}
\newtheorem{assumption}[theorem]{Assumption}
\newtheorem{observation}[theorem]{Observation}
\newtheorem*{remark}{Remark}
\newtheorem{openproblem}[theorem]{Open Problem}

% -- Macros 
\newcommand{\Hstr}{H_{\mathrm{str}}}
\newcommand{\Hnull}{H_{\mathrm{null}}}
\newcommand{\Estr}{E_{\mathrm{str}}}
\newcommand{\Espec}{E_{\mathrm{spec}}}
\newcommand{\etaorig}{\eta_{\mathrm{orig}}}
\newcommand{\Tren}{T_{\mathrm{ren}}}
\newcommand{\Gun}{G^{\mathrm{un}}}
\newcommand{\Gnorm}{G^{\mathrm{norm}}}
\newcommand{\acan}{a_p^{\mathrm{can}}}
\newcommand{\Dcan}{D^{\mathrm{can}}}
\newcommand{\Estrcan}{\Estr^{\mathrm{can}}}
\newcommand{\Especcan}{\Espec^{\mathrm{can}}}
\newcommand{\Hlocal}{H_{\mathrm{local}}}
\newcommand{\Hren}{H_{\mathrm{ren}}}
\newcommand{\cren}{c_p^{\mathrm{ren}}}
\newcommand{\ceta}{c_p^{\eta}}
\newcommand{\fp}{f_p}

% -- Title 
\title{\vspace{-1.5em}\bfseries A Finite-Cutoff Hilbert-Space Model\\[0.4em]
\large for Prime--Zero Energy Structure}
\author{Ulrich Tehrani}
\date{Research Note \textperiodcentered\ August~2026 \textperiodcentered\ v14}

\begin{document}
\maketitle

\begin{abstract}
We construct a finite-cutoff Hilbert-space model for the spectral structure
associated with primes and nontrivial zeros of the Riemann zeta function.
Given a prime cutoff $\kappa$ and truncated zero ordinates
$\{\gamma_1,\ldots,\gamma_N\}$, we define two finite-dimensional real
Hilbert spaces $\Hstr$ (over primes) and $\Hnull$ (over zeros),
connected by a linear map $\Phi\colon\Hstr\to\Hnull$, and the
self-adjoint loop operator $T = \Phi^*\circ\Phi$.

\textbf{What is proved.}
We establish the algebraic identity
$\Delta := \Estr - \Espec
 = \sum_p c_p^2\|a_p\|^2 - \|\sum_p c_p a_p\|^2$
by direct expansion.

\textbf{What is numerical.}
$\Delta > 0$ and
$\etaorig = \Delta/\Estr \in [0.650,\,0.700]$ for benchmark parameters
$(\kappa,\varepsilon,N)=(53,0.05,100)$, $\sigma\in[0.1,0.95]$.

\textbf{What remains open.}
An analytic proof of $\etaorig>0$ remains open.

\textbf{Structure.}
The formal core in \S\S\,2--4 is non-circular and proof-complete;
results from Papers~1--2 appear only as context.
No use of the Riemann Hypothesis, the GUE conjecture,
or the Hilbert--P\'olya postulate is made anywhere in \S\S\,2--4.
The model suggests a picture in which $\sigma=\tfrac12$ plays
the role of an energy reference point.
\end{abstract}

\vspace{2em}

% --------------------------------------------------------
% Notation and Conventions
% --------------------------------------------------------
\section*{Notation and Conventions}
\label{sec:notation}

\emph{Critical line.}
Throughout this series, $\sigma$ denotes $\Re(s)$ and the critical
line means $\sigma=\tfrac12$. The functional equation symmetry
$\xi(s)=\xi(1-s)$ makes $\sigma=\tfrac12$ the natural reference
point.

\medskip
\emph{Global parameters.}
In this paper the reference benchmark values are
$\kappa=53$ (16~active primes $p\leq\kappa$),
$\varepsilon=0.05$, and $N=100$.
The formal definitions allow arbitrary fixed $\kappa>1$,
$\varepsilon>0$, and $N\geq 1$; numerical sections state
explicitly when these parameters are varied.

\medskip
\emph{Main objects.}
The prime space $\Hstr=\ell^2(S(\kappa))$ and zero space
$\Hnull=\ell^2(\{\gamma_1,\ldots,\gamma_N\})$ are connected by
$\Phi\colon\Hstr\to\Hnull$, $\Phi(e_p)=a_p$, where
$(a_p)_k=e^{-\varepsilon^2\gamma_k^2/2}\sin(\gamma_k\log p)$.
The loop operator is $T=\Phi^*\circ\Phi$\footnote{We call $T=\Phi^*\Phi$
the \emph{loop operator}, following the notation of this series. It is the
Gram matrix of the phasor family $\{a_p\}_{p\le\kappa}$ in
$H_{\mathrm{null}}$, and should not be confused with loop operators in
lattice gauge theory or condensed-matter physics.} and
$\etaorig=1-\Espec/\Estr$ measures the energy gap.

\medskip
\emph{Phasor convention.}
Throughout this paper the prime phasor is
$\sin(\gamma_k\log p)$ at $\sigma=\tfrac12$;
no additional $\sigma$-factor enters the construction.
This convention is specific to Papers~1--3 of this series.

\medskip
\emph{Weights.}
The local prime curvature at the critical line is
\[
  V_p\!\bigl(\tfrac12\bigr) \;=\; \frac{4(\log p)^2\,p}{(p-1)^2}.
\]
For reference, the renormalized weight of Paper~2 is
$f_p := V_p(\tfrac12) - 4(\log p)^2/p = 4(\log p)^2(2p-1)/(p(p-1)^2)$.
The canonical weight in the $\eta$-framework of this paper is
\[
  \ceta \;:=\; \sqrt{V_p\!\bigl(\tfrac12\bigr)}\,.
\]
The energy-asymmetry functional is defined as
\[
  \eta(c) \;:=\; 1 - \frac{c^\top G^{\mathrm{un}} c}{E_{\mathrm{str}}(c)}.
\]
In this paper $c = \ceta = \sqrt{V_p(\tfrac12)}$,
and we write $\etaorig := \eta(\ceta)$.
This specialization is distinct from Paper~4's $\eta_{\mathrm{ren}} := \eta(c^{\mathrm{ren}})$,
which uses $c^{\mathrm{ren}} = \sqrt{f_p}$ and gives a numerically different value.

\medskip
\emph{Not to be confused with.}
$T=\Phi^*\Phi$ (loop operator on $\Hstr$, this paper) is distinct from
$\tilde T=\Phi\Phi^*$ (the adjoint loop operator on $\Hnull$).
$D\approx 9.470$ (renormalized energy sum, Paper~2) is distinct from
$\Estr$ (structural energy at benchmark parameters, this paper).
The canonical weight $\ceta = \sqrt{V_p(\tfrac12)}$ in the
$\eta$-framework of this paper is \emph{distinct} from
Paper~2's renormalized weight $\cren = \sqrt{\fp}$:
both in formula ($V_p(\tfrac12) \neq \fp$) and in role
($\ceta$ enters $\Estr$ inside the $\eta$-formula, whereas
$\cren$ enters $D=\sum_p (\cren)^2$ in the Weil renormalization).

\medskip
\emph{Analytical foundations.}
The construction uses real finite-dimensional inner-product space theory
(Gram matrices, positive semi-definite operators) and the prime number theorem
via Mertens' formula for weight asymptotics. The explicit formula and Weil's
quadratic form~\cite{Weil} provide the infinite-dimensional framework from
which the finite-cutoff model is abstracted.

\medskip
\emph{Numerical computations.}
All numerical results were computed with \texttt{mpmath} at 50~decimal digits
of working precision; Riemann zero ordinates $\gamma_k$ were evaluated via
\texttt{mpmath.zetazero} and \texttt{numpy} was used for matrix operations.
Verification code is available at
\url{https://github.com/utehrani/analysislab-nt}.

\newpage

%% ═══════════════════════════════════════════════════════════════════════════
\section{Introduction}
%% ═══════════════════════════════════════════════════════════════════════════

\subsection{Background and Motivation}

The present paper is the third in a series developing an
operator-theoretic framework for studying the distribution of the
non-trivial zeros of the Riemann zeta function $\zeta(s)$.
The formal core in \S\S\,2--4 is non-circular and proof-complete.

\textbf{Paper~1}~\cite{Paper1} introduced the curvature function
$H_{\xi}(\sigma,t) := \frac{d^2}{d\sigma^2}\log|\xi(\sigma+it)|^2$
and established the divergence
$\Hlocal(\tfrac12,\kappa)\sim 2(\log\kappa)^2\to\infty$ at
$\sigma=\tfrac12$, contrasted with convergence for $\sigma>\tfrac12$.

\textbf{Paper~2}~\cite{Paper2} connected this local curvature to the
Weil explicit functional, constructing admissible test functions
and establishing the convergent weight sum
$D=\sum_p(c_p^{\mathrm{ren}})^2\approx 9.470$.
These results appear only as motivation in the introduction and in the
contextual remark of~\S\,2; the formal core does not depend on them.

\textbf{The present paper} makes the geometric structure underlying
both papers explicit: two finite-dimensional Hilbert spaces $\Hstr$
and $\Hnull$, connected by a linear map $\Phi$, and a self-adjoint
loop operator $T=\Phi^*\circ\Phi$.

Several established programs seek a spectral interpretation of the
Riemann zeros. Berry and Keating~\cite{BerryKeating} studied a
semiclassical analogue via quantum chaos;
Connes~\cite{Connes} realizes the zeros through a trace formula in
noncommutative geometry;
de~Branges~\cite{deBranges} approached the problem through Hilbert
spaces of entire functions;
Burnol~\cite{BurnolForum04} constructs a conductor operator on
Sonin spaces directly from the explicit formula.
These programs work in spectral or infinite-dimensional analytic
frameworks. The present model is structurally distinct:
$H_{\mathrm{str}}$ and $H_{\mathrm{null}}$ are finite-dimensional
and real, no spectral postulate is made, and the $\gamma_k$ enter
as labels of basis vectors, not as eigenvalues.
For a taxonomy of physics-of-RH approaches see
Schumayer--Hutchinson~\cite{SchumayerHutchinson11}.

The Gram-matrix construction $T_{pq}=\langle a_p,a_q\rangle$ from
damped sine vectors appears to be new; conceptually it may be viewed
as a discrete analogue of spectral interpretations of the explicit
formula, as studied by Goldfeld~\cite{Goldfeld}. The test functions of
Paper~2 connect to adelic Schwartz functions as in Tate's
thesis~\cite{Tate}; the Weil positivity criterion~\cite{Weil,Lagarias}
provides the outer framework.
Bombieri--Lagarias~\cite{BombieriLagarias99} establish that
finite-truncation positivity arguments are not zeta-specific;
the arithmetic content of this paper lies entirely in the numerical
value of $\etaorig>0$, not in the algebraic identity $\Delta\geq0$.
Bender--Brody--M\"uller~\cite{BenderBrodyMuller17} (with
rebuttal~\cite{Bellissard17}) propose a non-Hermitian Hamiltonian
whose reality of spectrum would imply $\sigma=\tfrac12$; our
approach treats $\sigma=\tfrac12$ as a finite-cutoff variance
reference point, motivated by functional-equation symmetry and
the local-curvature divergence from \textbf{Paper~1}~\cite{Paper1},
rather than spectral assumptions.
Ionescu~\cite{Ionescu19} forms a related prime-character pairing
without the Gaussian regulator or the bipartite operator framing.
Bombieri~\cite{Bombieri2000} gives a systematic analysis of the
Weil quadratic functional in the context of obstructions to RH;
our finite-cutoff model implements a truncated version of this
functional without addressing the infinite-dimensional obstructions.

The renormalized framework of Paper~2 and the $\eta$-framework of
this paper use different normalizations; a precise identification
remains open (Open Problem~\ref{op:shell}).

\subsection{Main Results}

We define two finite-dimensional real Hilbert spaces $\Hstr$ (over
primes $p\le\kappa$) and $\Hnull$ (over zero ordinates
$\gamma_1,\ldots,\gamma_N$), connected by the linear map
$\Phi:\Hstr\to\Hnull$ built from Gaussian-damped prime resonance
vectors (\S\,3). The self-adjoint loop operator $T=\Phi^*\circ\Phi$
is the Gram matrix of these vectors (\S\,3).

We prove the algebraic identity $\Delta:=\Estr-\Espec$ by direct
expansion (Theorem~4.3). Numerically, $\etaorig=\Delta/\Estr\in
[0.650,0.700]$ for benchmark parameters
$(\kappa,\varepsilon,N)=(53,0.05,100)$, $\sigma\in[0.1,0.95]$
(Observation~\ref{obs:nonneg}). An analytic proof of $\etaorig>0$ remains open
(Open Problem~\ref{op:pos}).

\subsection{Reader Contract}

This paper \emph{proves}, by direct algebraic expansion, the
variance identity $\Delta=\Estr-\Espec$ (Theorem~4.3) and the basic
properties of the loop operator $T=\Phi^*\circ\Phi$: self-adjointness,
positive semi-definiteness, and $\sigma$-invariance of the normalized
Gram matrix (Proposition~3.5). These are unconditional.

This paper \emph{establishes numerically} that $\etaorig>0$ for all
tested parameter combinations, with
$\etaorig\in[0.650,0.700]$ at benchmark parameters (\S\,4).
The finite-grid spectral bound $B^E_{\max}<1$ and $B^{\|c\|}_{\max}<1$ and the
analytic explanation of the value of $B^{\|c\|}_{\max}$ (Problem~5.4)
are numerical observations and open problems respectively.

This paper \emph{leaves open} the analytic proof of $\etaorig>0$
(Problem~5.1), the spectral bound for canonical weights
(Problem~5.2), the analytical cancellation bound (Problem~5.3),
the analytic explanation of $B^{\|c\|}_{\max}$ (Problem~5.4), and the
renormalized shell energy (Problem~5.5).

%% ═══════════════════════════════════════════════════════════════════════════
\section{Two Hilbert Spaces}
%% ═══════════════════════════════════════════════════════════════════════════

We work throughout with real Hilbert spaces; all vectors and operators have
real entries. Complex phasor language appears in \S\,3.1 as
heuristic motivation only; all formal objects ($\Hstr, \Hnull, \Phi, T$)
are real.

Throughout this paper, $\kappa>1$ is a fixed real cutoff parameter,
$S(\kappa):=\{p\text{ prime}:p\leq\kappa\}$ the set of active primes,
and $\{\gamma_k\}_{k=1}^N$ the first $N$ positive imaginary parts of
nontrivial zeros of $\zeta(s)$, ordered $0<\gamma_1\leq\gamma_2\leq\cdots$.
The Gaussian damping parameter $\varepsilon>0$ is fixed throughout.

\begin{definition}[Prime space $\Hstr$]
\[
\Hstr := \ell^2(S(\kappa))
\]
with basis $\{e_p : p\in S(\kappa)\}$, standard inner product
$\langle a,b\rangle_{\mathrm{str}}=\sum_p a_p b_p$, and
$\dim\Hstr=|S(\kappa)|=\pi(\kappa)$.
The \emph{structural energy} of a weight vector $c\in\Hstr$ is
$\Estr(c)=\sum_p c_p^2\|a_p\|^2$ (defined in \S4 after $\Phi$ is introduced).
\end{definition}

\begin{definition}[Zero space $\Hnull$]
\[
\Hnull := \ell^2(\{\gamma_1,\ldots,\gamma_N\})
\]
with basis $\{f_k:k=1,\ldots,N\}$, inner product
$\langle u,v\rangle_{\mathrm{null}}=\sum_{k=1}^N u_k v_k$
(unweighted), and $\dim\Hnull=N$.
\end{definition}

\begin{remark}[Why $\sigma=\tfrac12$ is a distinguished reference point]
\label{rem:half}
This remark provides context; it is not used in \S\S\,3--5.

Two independent proven reasons support this choice:

\emph{(i) Symmetry} [proven]: $\sigma=\tfrac12$ is the unique fixed point of
$s\mapsto 1-s$, the symmetry of the completed zeta function
$\xi(s)=\xi(1-s)$.

\emph{(ii) Curvature singularity} [proven, Papers~1--2]:
$\Hlocal(\tfrac12,\kappa)\sim 2(\log\kappa)^2\to\infty$ as $\kappa\to\infty$,
while $\Hlocal(\sigma,\kappa)\to C(\sigma)<\infty$ for all $\sigma>\tfrac12$.

$\sigma=\tfrac12$ is thus a distinguished reference point by symmetry and
curvature divergence. Energy near-equality ($\Estr\approx\Espec$) is NOT
observed in the finite-cutoff model (see \S5).
\end{remark}

%% ═══════════════════════════════════════════════════════════════════════════
\section{The Linear Map $\Phi$ and the Loop Operator $T$}
%% ═══════════════════════════════════════════════════════════════════════════

\emph{Convention.} Gaussian damping is carried by the vectors:
$a_p=(e^{-\varepsilon^2\gamma_k^2/2}\sin(\gamma_k\log p))_k$.
The $\Hnull$ inner product is unweighted: $\langle u,v\rangle=\sum_k u_kv_k$.

\subsection{Motivation for the resonance vectors}

The formal definition is Definition~3.2 below; this subsection
provides intuition.

The prime resonance vector at zero ordinate $\gamma_k$ and prime $p$ arises as
the imaginary part of the prime phasor $\varphi_p(u)=e^{iu\log p}$ evaluated
at $u=\gamma_k$, with Gaussian regularization for convergence:
\[
(a_p)_k = e^{-\varepsilon^2\gamma_k^2/2}\cdot\sin(\gamma_k\log p).
\]
The $k$-th component measures how strongly prime $p$ resonates with zero
ordinate $\gamma_k$.

\begin{definition}[Prime resonance vectors]
For each prime $p\in S(\kappa)$, the \emph{canonical resonance vector}
at $\sigma=\tfrac12$ is
\[
\acan := a_p :=
\bigl(e^{-\varepsilon^2\gamma_k^2/2}\sin(\gamma_k\log p)\bigr)_{k=1}^N
\;\in\;\Hnull.
\]
At general $\sigma$, the resonance vectors depend on $\sigma$ via
\[
a_p^{(\sigma)} := p^{-(\sigma-1/2)}\cdot\acan.
\]
The vectors $a_p=\acan=a_p^{(1/2)}$ at $\sigma=\tfrac12$ are the
canonical representatives used in all primary definitions.
\end{definition}

\begin{definition}[Linear map $\Phi$]
\[
\Phi\colon\Hstr\to\Hnull,\qquad\Phi(e_p):=\acan.
\]
Extended linearly: $\Phi(c)=\sum_p c_p\acan$ for any
$c=\sum_p c_p e_p\in\Hstr$.

\emph{We define $\Phi$ as a linear map; no injectivity claim is made for the
finite-cutoff model.}
\end{definition}

\begin{definition}[Gram matrices]
The \emph{unnormalized Gram matrix} (canonical) is
\[
\Gun_{pq} := \Gun{}^{\mathrm{,can}}_{pq}
:= \langle\acan,a_q^{\mathrm{can}}\rangle_{\Hnull}
= \sum_{k=1}^N e^{-\varepsilon^2\gamma_k^2}\sin(\gamma_k\log p)\sin(\gamma_k\log q).
\]
The $\sigma$-dependent extension is
$\Gun(\sigma)_{pq}=p^{-(\sigma-1/2)}q^{-(\sigma-1/2)}\Gun_{pq}$.
The \emph{normalized Gram matrix} is
\[
\Gnorm_{pq}:=\frac{\Gun_{pq}}{\|\acan\|\cdot\|a_q^{\mathrm{can}}\|}
=\Bigl\langle\frac{\acan}{\|\acan\|},\frac{a_q^{\mathrm{can}}}{\|a_q^{\mathrm{can}}\|}\Bigr\rangle_{\Hnull}.
\]
The matrix $\Gnorm$ is $\sigma$-invariant (Proposition~3.5(iv)).
\end{definition}

\begin{definition}[Loop operator $T$]
\[
T:=\Phi^*\circ\Phi\colon\Hstr\longrightarrow\Hstr.
\]
The matrix representation of $T$ in the prime basis is $T_{pq}=\Gun_{pq}$.
$T$ is thus identical to the canonical Gram matrix $\Gun$, and is explicitly
constructed from prime resonance vectors.
\end{definition}

\begin{proposition}[Properties of $T$; proved]
\label{prop:T}
\begin{enumerate}[label=\emph{(\roman*)}, leftmargin=2.2em]
\item \emph{Self-adjoint:} $T^*=(\Phi^*\Phi)^*=\Phi^*\Phi=T$.
\item \emph{Positive semi-definite:} $\langle Tc,c\rangle=\|\Phi c\|^2\geq 0$
for all $c\in\Hstr$.
\item \emph{Non-negative eigenvalues:} All eigenvalues $\lambda_j\geq 0$.
\item \emph{$\sigma$-invariance of $\Gnorm$} [by direct calculation]:
For the $\sigma$-dependent family $a_p^{(\sigma)}$, the normalized Gram matrix
$\Gnorm_{pq}$ is independent of $\sigma$. \emph{Proof.} The scalar factor
$p^{-(\sigma-1/2)}$ cancels in the normalized ratio. \qed
\end{enumerate}
\end{proposition}

\begin{remark}[Spectral structure of $T$]
The eigenvalues $\lambda_j\geq 0$ of $T=\Phi^*\circ\Phi$ are the squares of
the singular values $\sigma_j(\Phi)$: $\lambda_j=\sigma_j(\Phi)^2$.
The eigenvalues $\lambda_j$ are built from the zero ordinates $\gamma_k$
appearing in the definition of $\acan$, but are \emph{not} directly the zeros
$\gamma_k$ themselves. This distinguishes $T$ from the classical Hilbert-P\'olya
ansatz: $T$ is explicitly constructed, not postulated.
\end{remark}

%% ═══════════════════════════════════════════════════════════════════════════
\section{The Variance Identity and $\etaorig$}
%% ═══════════════════════════════════════════════════════════════════════════

\emph{Convention.} Objects without $\sigma$-argument refer to the canonical
($\sigma=\tfrac12$) version throughout \S\,4.

\begin{definition}[Canonical and $\sigma$-dependent energies]
\label{def:energies}
For a weight vector $c\in\Hstr$ with components $c_p$:

\emph{Canonical quantities} (at $\sigma=\tfrac12$):
\[
\Dcan:=\mathrm{diag}\!\bigl(\|\acan\|^2\bigr)_{p\in S(\kappa)},
\qquad
\Estrcan:=c^\top\Dcan c=\sum_p c_p^2\|\acan\|^2,
\]
\[
\Especcan:=\|\Phi c\|^2=c^\top\Gun c.
\]

\emph{$\sigma$-dependent extensions:}
\[
D(\sigma)_{pp}:=p^{-2(\sigma-1/2)}\cdot\Dcan_{pp},
\qquad
\Estr(\sigma):=\sum_p c_p^2\,p^{-2(\sigma-1/2)}\|\acan\|^2,
\]
\[
\Espec(\sigma):=c^\top\Gun(\sigma)\,c.
\]

The \emph{canonical weight vector} in the $\eta$-context is
\[
\ceta \;:=\; \sqrt{V_p\!\bigl(\tfrac12\bigr)}
    \;=\; \sqrt{\frac{4(\log p)^2\,p}{(p-1)^2}},
    \qquad p \in S(\kappa).
\]
\begin{remark}[Canonical weights --- normative form]
The canonical weight $\ceta = \sqrt{V_p(\tfrac12)}$
is the normative form in the $\eta$-framework. For large $p$,
$\ceta \approx 2(\log p)/\!\sqrt{p}$; this differs from the Weil weight
$\cren = \sqrt{\fp}$ both in formula and asymptotic behaviour.
\end{remark}
\end{definition}

\begin{assumption}[Positive definiteness of $D$]
Throughout this paper we assume $\Dcan_{pp}=\|a_p\|^2>0$ for all
$p\in S(\kappa)$. This holds in all numerical experiments conducted here.
\end{assumption}

\begin{theorem}[Algebraic identity; proved]
\label{thm:identity}
The following equality holds by direct expansion:
\[
\Delta:=\Estr-\Espec=\sum_p c_p^2\|a_p\|^2-\Bigl\|\sum_p c_p a_p\Bigr\|^2.
\]
\end{theorem}

\begin{proof}
Expand $\Espec=\|\Phi c\|^2=\bigl\langle\sum_p c_p a_p,\sum_q c_q
a_q\bigr\rangle=\sum_{p,q}c_p c_q\langle a_p,a_q\rangle=c^\top\Gun c$.
Then $\Delta=c^\top\Dcan c-c^\top\Gun c$, which is the stated identity.
\end{proof}

\begin{remark}[Algebraic identity vs.\ arithmetic content]
The inequality $\Delta\geq0$ is \emph{not} an algebraic consequence of
Theorem~\ref{thm:identity}: $T=\Phi^*\Phi$ is positive semi-definite,
but $\Delta=c^\top(\Dcan-\Gun)c$ has indefinite sign for a general
weight vector $c$. The non-trivial arithmetic content is the numerical
lower bound $\etaorig\geq0.650>0$ for the canonical weight
$\ceta=\sqrt{V_p(\tfrac12)}$ at the reference parameters, established
in the observations of this section.
\end{remark}

\begin{observation}[Non-negativity, numerical]
\label{obs:nonneg}
\[
\Delta\geq 0\quad\text{equivalently}\quad\etaorig\geq 0
\]
holds for all tested parameters with canonical weights
$\ceta=\sqrt{V_p(\tfrac12)}$ (Definition~\ref{def:energies}).
\emph{[Numerical only --- no unconditional proof for general $c$.]}

Note: For \emph{all} weight vectors $c$, the inequality $\etaorig(c)\geq 0$
would be equivalent to
$\lambda_{\max}\bigl((\Dcan)^{-1/2}T(\Dcan)^{-1/2}\bigr)\leq 1$.
This universal bound is false on the tested cutoffs: on those cutoffs,
$\lambda_{\max}(\Tren)/\pi(\kappa)$ is numerically of order $0.39$, so
$\lambda_{\max}(\Tren)>1$ already at finite tested cutoffs.
The relevant open problem is therefore not a universal bound in $c$,
but the canonical-weight Rayleigh quotient condition in
Open Problem~\ref{op:spectral}.

\emph{Cauchy-Schwarz does \textbf{not} guarantee $\etaorig\geq 0$}: the
correct identity is $\Delta=c^\top(\Dcan-\Gun)c$; the sign of $\Delta$ is not
determined by the algebraic identity alone.
\end{observation}

\begin{definition}[$\etaorig$]
\label{def:eta}
\[
\etaorig:=1-\frac{\Espec}{\Estr}=1-\frac{\|\Phi c\|^2}{\Estr}.
\]
Boundary cases: $\etaorig=0$ iff $\Espec=\Estr$; $\etaorig=1$ iff $\Phi c=0$.
\end{definition}

\begin{remark}[$\sigma$-dependence via rescaled weights]
The entire $\sigma$-dependence of $\etaorig(\sigma)$ is carried by the
rescaled weight vector $\tilde{c}_p(\sigma):=c_p\cdot p^{-(\sigma-1/2)}$:
\[
\etaorig(\sigma)=1-\frac{\tilde{c}(\sigma)^\top\Gun\,\tilde{c}(\sigma)}
{\tilde{c}(\sigma)^\top\Dcan\,\tilde{c}(\sigma)}.
\]
\emph{Proof.} Since $\Gun(\sigma)_{pq}=p^{-(\sigma-1/2)}q^{-(\sigma-1/2)}\Gun_{pq}$,
one has $c^\top\Gun(\sigma)c=\tilde{c}^\top\Gun\tilde{c}$ and
$\Estr(\sigma)=\tilde{c}^\top\Dcan\tilde{c}$. \qed
\end{remark}

\emph{Numerical evidence.}
All computations use $N=100$ zero ordinates at the reference
parameters; the algebraic identity (Theorem~\ref{thm:identity}) is
verified to rounding error $<10^{-14}$. Verification scripts are
available in the companion repository~\cite{Code}.

\begin{observation}[$\etaorig$ positive across all tested parameters, numerical]
$\etaorig>0$ in all tested parameter combinations. For the benchmark slice
$(\kappa,\varepsilon,N)=(53,0.05,100)$, $\sigma\in[0.1,0.95]$:
\[
\etaorig\in[0.650,\;0.700].
\]
Note: the range $[0.650,0.700]$ is strongly $\varepsilon$-dependent; no
universal range over all $(\kappa,\varepsilon)$ is claimed.
\end{observation}

The positivity check was performed for
$\kappa\in\{23,53,101,199\}$, $\varepsilon=0.05$, $N=100$,
and the $\sigma$-grid $\{0.10,0.20,0.30,0.40,0.50,0.60,0.70,0.80,0.90,0.95\}$.
The table below displays the benchmark slice
$(\kappa,\varepsilon,N)=(53,0.05,100)$: on this grid $\etaorig(\sigma)$ has
a shallow local maximum near $\sigma\approx0.20$ and then decreases.
No monotonicity theorem is claimed.

\begin{center}
\begin{tabular}{cc}
\toprule
$\sigma$ & $\etaorig$ \\
\midrule
0.10 & 0.69928647 \\
0.20 & 0.69989546 \\
0.30 & 0.69866521 \\
0.40 & 0.69558817 \\
0.50 & 0.69078176 \\
0.60 & 0.68441702 \\
0.70 & 0.67662272 \\
0.80 & 0.66742445 \\
0.90 & 0.65675975 \\
0.95 & 0.65085191 \\
\bottomrule
\end{tabular}
\end{center}

\begin{observation}[Finite-grid mode diagnostics $B^E_{\max}$ and $B^{\|c\|}_{\max}<1$, numerical]
\label{obs:bmax}
The following computations use
$\kappa\in\{23,53,101,199\}$, $\varepsilon\in\{0.02,0.05,0.10\}$,
$\sigma\in\{0.3,0.4,0.5,0.6,0.7\}$, $N=100$.

Set $\tilde{c}(\sigma)_p:=c_p\cdot p^{-(\sigma-1/2)}$ (rescaled weight vector).
Two mode-contribution diagnostics are computed, where $\lambda_j,u_j$
denote eigenvalues and eigenvectors of $T=\Gun$ (at $\sigma=\tfrac12$,
canonical):
\[
B^E_j(\sigma):=\frac{\lambda_j\,|\langle\tilde{c}(\sigma),u_j\rangle|^2}
               {\tilde{c}(\sigma)^\top D^{\mathrm{can}}\tilde{c}(\sigma)},
\qquad
B^{\|c\|}_j:=\frac{\lambda_j\,|\langle c,u_j\rangle|^2}{\|c\|^2}.
\]
If no eigenvalue exceeds~1, we set $B_{\max}:=0$ by convention.
The $E_{\mathrm{str}}$-normalized diagnostic satisfies
$\sum_j B^E_j(\sigma)=1-\etaorig(\sigma)$ [algebraic, all $\sigma$];
the $\|c\|^2$-normalized diagnostic does not.
Numerically at $(\kappa,\varepsilon)=(53,0.05)$:
\begin{align*}
B^E_{\max}(\sigma) &:= \max_{j:\lambda_j>1} B^E_j(\sigma)
\in[0.115,\,0.127]\text{ for }\sigma\in[0.3,0.7],\quad <1 \\
B^{\|c\|}_{\max} &:= \max_{j:\lambda_j>1} B^{\|c\|}_j
\approx 0.0735,\quad<1
\end{align*}
Both satisfy their respective maximum $<1$ for all tested parameters.
\end{observation}

Selected values of $B^{\|c\|}_{\max}$ (the tabulated diagnostic):

\begin{center}
\begin{tabular}{cccc}
\toprule
$\kappa$ & $\varepsilon=0.02$ & $\varepsilon=0.05$ & $\varepsilon=0.10$ \\
\midrule
23  & 0.2975 & 0.0112 & 0.0000 \\
53  & 0.3813 & 0.0735 & 0.0000 \\
101 & 0.1241 & 0.0483 & 0.0000 \\
199 & 0.1546 & 0.0483 & 0.0107 \\
\bottomrule
\end{tabular}
\end{center}

\begin{observation}[$B^{\|c\|}_{\max}$ $\sigma$-invariant by construction, finite-grid numerical]
The diagnostic $B^{\|c\|}_{\max}\approx0.0735$ for $(\kappa,\varepsilon)=(53,0.05)$
is independent of $\sigma$ by construction: both $c$ and the eigenvectors $u_j$
of $T$ are fixed at $\sigma=\tfrac12$. The analytical explanation of the
value $0.0735$ remains open (Open Problem~\ref{op:sigma}).
\end{observation}

\begin{remark}[Cancellation mechanism, numerical analysis]
The mechanism underlying $B_{\max}<1$ is a cancellation between the monotone

weight vector $\ceta=\sqrt{V_p(\tfrac12)}$ and the oscillatory eigenvectors $u_j$
of $\Gun$: for $\kappa=53$, eigenvectors $u_1,\ldots,u_4$ each exhibit 10
sign changes along the prime ordering. Cancellation rates 67--97\% observed
for $\kappa\in\{23,53,101,199\}$. Formal proof: Open Problem~\ref{op:cancel}.
\end{remark}

%% ═══════════════════════════════════════════════════════════════════════════
\section{Open Problems}
%% ═══════════════════════════════════════════════════════════════════════════

Unless otherwise stated, the open problems below concern the finite-cutoff
family with $T=T(\kappa,\varepsilon,N)$,
$D^{\mathrm{can}}=D^{\mathrm{can}}(\kappa,\varepsilon,N)$, and canonical
weights $c_p^\eta=\sqrt{V_p(\tfrac12)}$.
Numerical evidence uses $N=100$ and the explicitly stated grids in
$(\kappa,\varepsilon)$.
A uniform theorem in $\kappa$, $\varepsilon$, or an asymptotic regime is a
stronger statement and is not claimed unless specified.


\begin{openproblem}[Analytic positivity]\label{op:pos}
Prove analytically that $\etaorig>0$ for the canonical weight vector
$\ceta=\sqrt{V_p(\tfrac12)}$.
Numerically confirmed: $\etaorig\in[0.650,0.700]$ for $(\kappa,\varepsilon,N)=(53,0.05,100)$,
$\sigma\in[0.1,0.95]$ (\S\,4).
\end{openproblem}

\begin{openproblem}[Spectral bound for canonical weights]\label{op:spectral}
Let
\[
\Tren:=(\Dcan)^{-1/2}T(\Dcan)^{-1/2},\qquad x:=(\Dcan)^{1/2}c,
\]
where $c=c^\eta$ is the canonical weight vector. For this fixed canonical
weight, $\etaorig>0$ is equivalent to
\[
\frac{\langle\Tren x,x\rangle}{\|x\|^2}<1.
\]
The universal bound $\lambda_{\max}(\Tren)<1$, which would imply positivity
for every weight vector, is not true on the tested cutoffs and is not
claimed here.
Prove the above Rayleigh quotient condition analytically for
$c_p^\eta=\sqrt{V_p(\tfrac12)}$.
For all tested $\kappa\in\{23,53,101,199,503,1009\}$ the condition holds
numerically. The vector $x$ appears nearly orthogonal to the dominant
eigenmode of $\Tren$ (localized at $\gamma_1=14.13\ldots$); proving this
near-orthogonality is the core open problem.
\end{openproblem}

\begin{openproblem}[Analytical cancellation bound]\label{op:cancel}
Prove analytically that
\[
B^{\|c\|}_{\max} :=
\max_{j:\,\lambda_j>1}\lambda_j\,
\frac{|\langle c,u_j\rangle|^2}{\|c\|^2} < 1,
\]
or equivalently, for all $j$ with $\lambda_j>1$:
\[
\frac{|\langle c,u_j\rangle|^2}{\|c\|^2}\leq\frac{C_B}{\lambda_j},\quad C_B<1.
\]
Note: since the tested data show growth of $\lambda_{\max}(\Tren)$
with $\pi(\kappa)$,
a uniform projection bound $|\langle c,u_j\rangle|^2/\|c\|^2\leq C_1\ll 1$
does not in itself imply $B^{\|c\|}_{\max}<1$; the growth of $\lambda_j$
must also be absorbed.
Approach: Abel summation for $c_p^\eta=\sqrt{V_p(\tfrac12)}\sim 2(\log p)/\!\sqrt{p}$
against oscillatory partial sums $\sum_{q\leq p}(u_j)_q$, without reference to RH.
Numerically: $B^{\|c\|}_{\max}<0.382$ for the tested grid,
$\kappa\in\{23,53,101,199\}$, $\varepsilon\in\{0.02,0.05,0.10\}$, $N=100$.
\end{openproblem}

\begin{openproblem}[Analytic explanation of $B^{\|c\|}_{\max}$]\label{op:sigma}
The diagnostic $B^{\|c\|}_{\max}:=\max_{j:\lambda_j>1}\lambda_j|\langle c,u_j\rangle|^2/\|c\|^2$
uses the fixed canonical weight $c=c_p^\eta$ and canonical eigenvectors $u_j$
of $T$ at $\sigma=\tfrac12$. It therefore carries no $\sigma$-dependence by
definition, and its $\sigma$-invariance is a structural feature of the formula.
What remains open is the analytical explanation of the numerical value
$B^{\|c\|}_{\max}\approx0.0735$ at $(\kappa,\varepsilon)=(53,0.05)$.
\end{openproblem}


\begin{openproblem}[Renormalized shell energy]\label{op:shell}
In the Paper~2 normalization framework, define the renormalized shell energy
using canonical effective coupling $N_{\mathrm{eff}}=\pi(\kappa)$:
\[
H_{\mathrm{ren}}(\kappa,\varepsilon(\kappa))
:= H_{\mathrm{local}}\!\bigl(\tfrac12,\kappa\bigr) - 2(\log\kappa)^2
= O(\log\kappa)
\]
(Paper~2, \S5), where $\pi(\kappa)$ is the prime-counting function
at cutoff $\kappa$.
Compute $H_{\mathrm{ren}}(\kappa,\varepsilon(\kappa))$ explicitly and
determine whether it has a stable limiting or asymptotic form.
This problem belongs to the Paper~2 normalization and is independent of
the algebraic identity proved in Section~4.
\end{openproblem}

\medskip
\emph{Scope and limitations.}
The algebraic identity
$\Delta = E_{\mathrm{str}} - E_{\mathrm{spec}}
= c^\top D^{\mathrm{can}}c - c^\top G^{\mathrm{un}}c$
is a direct consequence of the definitions. It does \emph{not}
imply $\Delta\ge0$: for a general weight vector $c$, the form
$D^{\mathrm{can}}-G^{\mathrm{un}}$ need not be positive
semidefinite. The sign statement $\Delta\ge0$ is established
here only numerically for the canonical weight
$c_p^\eta=\sqrt{V_p(\tfrac12)}$ on the tested parameter grids.
Bombieri--Lagarias~\cite{BombieriLagarias99} establish that
finite-truncation positivity is generically not zeta-specific;
the arithmetic content lies in the specific numerical value
$\etaorig\in[0.650,0.700]$ and in the five open problems.
No claim is made toward the Riemann Hypothesis.

%% ═══════════════════════════════════════════════════════════════════════════
\section*{Non-Circularity}
\label{sec:nocirc}
%% ═══════════════════════════════════════════════════════════════════════════

We document explicitly that the results of this paper do not rely on
the objects they are directed toward.

\textbf{No RH as input.}
At no point is the Riemann Hypothesis assumed.
The zero ordinates $\gamma_k$ enter the construction of $\Phi$ as
numerically specified inputs, and no property of their distribution
requiring the Riemann Hypothesis is assumed or used.

\textbf{No GUE, Montgomery, or Hilbert--P\'olya hypothesis.}
The loop operator $T=\Phi^*\Phi$ is built explicitly from the prime
resonance vectors $\{a_p\}_{p\le\kappa}$. Its construction postulates
neither GUE statistics, Montgomery's pair correlation conjecture, nor
the Hilbert--P\'olya conjecture. The eigenvalues $\lambda_j$ of $T$ are
the squares of the singular values of $\Phi$ and are not identified
with zero ordinates $\gamma_k$.

\textbf{Status separation maintained.}
The variance identity $\Delta=\Estr-\Espec$ is an algebraic theorem;
the sign statement $\Delta\geq 0$ is established numerically on the
tested grid. These are distinct levels and are flagged separately
throughout. No claim relating $\etaorig$ to RH is proved or used.

%% ═══════════════════════════════════════════════════════════════════════════
\section*{Acknowledgments}
%% ═══════════════════════════════════════════════════════════════════════════

The author thanks the anonymous readers who provided feedback on
earlier versions of this paper.

%% ═══════════════════════════════════════════════════════════════════════════
\section*{Call for Feedback}

This paper is part of an open research initiative on prime--zero
coupling and explicit-formula operators.
All papers in the series are freely available on Zenodo under CC~BY~4.0,
and the accompanying verification code is hosted on GitHub.

{\raggedright
Zenodo:\\
\url{https://zenodo.org/search?q=metadata.creators.person_or_org.name\%3A\%22Tehrani\%2C\%20Ulrich\%22}\\
GitHub:\\
\url{https://github.com/utehrani/analysislab-nt}
\par}

Topics of particular interest include: whether $\etaorig>0$ is an
expected consequence of known equidistribution results for the phases
$\{\gamma_k\log p \bmod 2\pi\}$, and whether the observed
near-orthogonality of the weight vector to the dominant eigenmode of
$T$ connects to standard cancellation phenomena in partial sums
involving primes and zeros.

%% ═══════════════════════════════════════════════════════════════════════════

\newpage
\begin{thebibliography}{99}

\bibitem{Weil}
A.~Weil,
``Sur les `formules explicites' de la th\'eorie des nombres premiers'',
\emph{Comm.\ S\'em.\ Math.\ Univ.\ Lund, Suppl.}\ (1952), 252--265.

\bibitem{Paper1}
U.~Tehrani,
``A Curvature Decomposition of the Explicit Formula
  for the Riemann Zeta Function'',
Zenodo, \url{https://doi.org/10.5281/zenodo.19025598}, 2026.

\bibitem{Paper2}
U.~Tehrani,
``From Local Curvature to the Weil Functional:
  An Explicit Construction'',
Zenodo, \url{https://doi.org/10.5281/zenodo.19106992}, 2026.

\bibitem{BerryKeating}
M.~V.~Berry and J.~P.~Keating,
``The Riemann zeros and eigenvalue asymptotics'',
\emph{SIAM Rev.} \textbf{41} (1999), 236--266.

\bibitem{Connes}
A.~Connes,
``Trace formula in noncommutative geometry and the zeros of the Riemann
zeta function'',
\emph{Selecta Math.\ (N.S.)} \textbf{5} (1999), 29--106.

\bibitem{deBranges}
L.~de~Branges,
``The Riemann hypothesis for Hilbert spaces of entire functions'',
\emph{Bull.\ Amer.\ Math.\ Soc.} \textbf{15} (1986), 1--17.

\bibitem{BurnolForum04}
J.-F.~Burnol,
``Two complete and minimal systems associated with the zeros of the
Riemann zeta function'',
\textit{Forum Math.}\ \textbf{16} (2004), 509--570;
arXiv:math/0103011.

\bibitem{SchumayerHutchinson11}
D.~Schumayer and D.~A.~W.~Hutchinson,
``Colloquium: Physics of the Riemann hypothesis'',
\textit{Rev.\ Mod.\ Phys.}\ \textbf{83} (2011), 307--330;
arXiv:1101.3116.

\bibitem{Goldfeld}
D.~Goldfeld,
``A spectral interpretation of Weil's explicit formula'',
in: J.~Jorgenson and S.~Lang (eds.),
\emph{Explicit Formulas for Regularized Products and Series},
Lecture Notes in Math.\ \textbf{1593}, Springer, 1994, pp.~136--152.
\url{https://doi.org/10.1007/BFb0074041}

\bibitem{Tate}
J.~T.~Tate,
``Fourier Analysis in Number Fields and Hecke's Zeta-Functions'',
in: J.~W.~S.~Cassels and A.~Fr\"ohlich (eds.),
\emph{Algebraic Number Theory},
Academic Press, 1967, pp.~305--347.

\bibitem{Lagarias}
J.C.~Lagarias,
``On a positivity property of the Riemann $\xi$-function'',
\emph{Acta Arithmetica} \textbf{89} (1999), 217--234.

\bibitem{BombieriLagarias99}
E.~Bombieri and J.~C.~Lagarias,
``Complements to Li's criterion for the Riemann hypothesis'',
\textit{J.\ Number Theory} \textbf{77} (1999), 274--287.

\bibitem{BenderBrodyMuller17}
C.~M.~Bender, D.~C.~Brody, and M.~P.~M\"uller,
``Hamiltonian for the zeros of the Riemann zeta function'',
\textit{Phys.\ Rev.\ Lett.}\ \textbf{118}, 130201 (2017).

\bibitem{Bellissard17}
J.~Bellissard,
``Comment on `Hamiltonian for the zeros of the Riemann zeta function'\,'',
arXiv:1704.02644 (2017).

\bibitem{Ionescu19}
L.~M.~Ionescu,
``A note on the statistics of Riemann zeros'',
arXiv:1903.09318 (2019).

\bibitem{Bombieri2000}
E.~Bombieri,
``Remarks on Weil's quadratic functional in the theory of prime
numbers,~I'',
\textit{Rend.\ Mat.\ Acc.\ Lincei}, Ser.~IX,
\textbf{11} (2000), 183--233.

\bibitem{Code}
U.~Tehrani,
Verification scripts for the prime--zero coupling series,
GitHub, \url{https://github.com/utehrani/analysislab-nt}.

\end{thebibliography}


% --------------------------------------------------------
% Footer
% --------------------------------------------------------
\enlargethispage{3\baselineskip}
\vspace*{\fill}
\begin{minipage}{\linewidth}
\rule{\linewidth}{0.4pt}\smallskip\\
\small\emph{Paper~3 v14 \textperiodcentered\ August~2026
\textperiodcentered\ Pauca sed matura}
\end{minipage}

\end{document}
